\scanpage{58}
\mspage{7}

$W$-fibré (au \uline{sens fort} a))
\[
S'\longrightarrow S_0\subset S
\]

avec $S'$ $W$-\uline{asphérique}. Comme $X\unsure{\text{---}|\longrightarrow}Y$ sur $S$ est dans $W$, il en sera donc de même de $X'\mathrel{\text{---}}Y'$ sur $S'$. \deleted{\unsure{Cela}} Comme les fibres de $X',Y'$ en $s'$ s'identifient à celles de $X,Y$ en l'image de $s'$ dans $S$, on est ramené au cas \inserted{où} $S$ \deleted{\unsure{asp}} est $W$-asphérique. Mais alors on a

\diagram{\begin{tikzcd}[ampersand replacement=\&,column sep=large]
X\arrow[r,"f"] \& Y\\
X_s\arrow[u,"i_s^X"]\arrow[r,"f_s"'] \& Y_s\arrow[u,"i_s^Y"']
\end{tikzcd}}

, où $i_s^X,i_s^Y$ sont dans $W$, car déduites de $\{s\}\longrightarrow S$ (qui est dans $W$, $S$ étant $W$-asphérique) par les chgt. de base $W$-fibrants $X\longrightarrow S$, $Y\longrightarrow S$. Donc $f\in W\Longrightarrow f_s\in W$.

NB On n'a utilisé la propriété de fibration pour $X,Y$ \unsure{au} sens a$'$).\textsuperscript{$*$}

% slanted
\margin{Et même, dans a$'$) on pourrait \deleted{\unsure{exiger}} que $Y''$ \unsure{soit} $=e$, $Y'$ $W_\omega$-asphérique, \uline{\unsure{page} 14}, $W_\omega\subset W\cap W_\infty$ ?}

Il faut donc \deleted{\unsure{voir}} prouver, si on peut, le

\uline{Lemme} (?) Soit $S$ un objet $0$-connexe de $\mathrm{Cat}$. Alors il existe un objet $W$-asphérique $S'$ de $\mathrm{Cat}$, et une \deleted{\unsure{$\mathrm{Hot}$}}-fibration \inserted{(au sens \uline{fort} de a) page 4)} $S'\longrightarrow S$.

% slanted
\margin{Est prouvé moyennant l'axiome \textcircled{4} ($+$ des petites vérifications \unsure{ultérieures})}

L'idée naturelle, c'est de prendre, pour $s_0\unsure{\subset|\in}\mathrm{Ob}\,S$ fixé,
\[
S'=\underline{\mathrm{Ch}}(S;s_0,-)\quad\text{ou}\quad S'=\underline{\mathrm{Ch}}_\infty(S;s_0,-),
\]
et la flèche
\[
\beta:c\longmapsto\mathrm{but}\,c:S'\longrightarrow S.
\]

On sait que $S'$ est $W_\infty$-asphérique et que $\beta$ est $W_\infty$-fibrante (et même parfaite). Il faut voir que $S'$ est aussi $W$ asphérique, et $\beta$ $W$-fibrante. Il faut donc revoir la démonstration dans le cas de $W_\infty$, et prouver qu'elle se \unsure{transfère|transforme} à $W$,

% archive sheet number: 58
